\chapter{Implementierung}
\label{ch:implementierung}

In Kapitel 4 wurde begr{\"u}ndet, was gemessen wird und warum. Zudem wurden in Kapitel 5 die einzelnen benchmark-Metriken vorgestellt. Dieses Kapitel beschreibt die softwaretechnische Umsetzung dieser Untersuchung. Der Fokus liegt dabei auf der Erkl{\"a}rung, wie die einzelnen Bauteile des Projekts zusammenarbeiten und welche Entwurfsentscheidungen jeweils getroffen wurden. Zuerst wird die grobe Organisation des Projekts vorgestellt (siehe 6.1). Anschlie{\ss}end wird die von allen Benchmarks gemeinsam genutzte Mess-Infrastruktur erl{\"a}utert (siehe 6.2). Zuletzt wird eine vollst{\"a}ndige Darstellung der umgesetzten Konfigurationsstufen gebildet.

\section{Projektstruktur}
Der West-Arbeitsbereich trennt zwischen zwei Bereichen. Bei dem ersten Bereich handelt es sich um einen \lstinline[style=fileinline]|zephyr|-Ordner, der die externe Zephyr-RTOS-Quelle enth{\"a}lt und von West selbst verwaltet wird und dem eigenen Projektordner, welcher den Namen \glqq{}thesis\grqq{} (\lstinline[style=fileinline]|thesis/|) tr{\"a}gt. Dieser enth{\"a}lt die gesamte, f{\"u}r diese Arbeit entwickelte Software und wird in diesem Kapitel beschrieben.

Die Datei \lstinline[style=fileinline]|thesis/west.yml| wird Manifest genannt. Diese legt fest, welche exakte Zephyr-Version (v3.7.2) und welche zugeh{\"o}rigen Module, darunter Hardware-Abstraktionsschicht, Bootloader, Netzwerk-Stack, etc. beim n{\"a}chsten west update in den \lstinline[style=fileinline]|zephyr|-Ordner geladen werden (siehe Kapitel 4.6.4). Dadurch ist die verwendete Software-Version nicht davon abh{\"a}ngig, welche Software-Version zuf{\"a}llig installiert wurde, sondern in einer einzigen, versionierten Datei festgelegt. Dies ist hierbei besonders f{\"u}r die Reproduzierbarkeit wichtig.

Eine kleine Konfigurationsdatei, \lstinline[style=fileinline]|thesis/zephyr/module.yml|, teilt dem Zephyr-Build-System mit, dass \lstinline[style=fileinline]|benchlib| als sogenanntes Modul behandelt werden soll. Ohne diese Datei w{\"u}rde Zephyr die eigene Bibliothek beim Bauen nicht finden. Diese Datei zeigt zudem nur auf zwei Pfade; auf die CMake-Bauvorschrift und auf die Kconfig-Definition der Bibliothek.

In Kapitel 6.2 wird im Detail die gemeinsame Mess-Infrastruktur erkl{\"a}rt. Diese ist unterteilt in:

\begin{itemize}
\item eine {\"o}ffentliche Schnittstelle (siehe 6.2.2): \lstinline[style=fileinline]|include/benchlib.h|
\item drei Implementierungsdateien (siehe 6.2.3 -- 6.2.5): \lstinline[style=fileinline]|src/|
\item eine CMake-Datei, welche festlegt, dass diese drei Dateien nur dann {\"u}bersetzt werden, wenn eine Anwendung \lstinline[style=fileinline]|CONFIG_BENCHLIB=y| setzt. Dadurch bleibt die Bibliothek folgenlos f{\"u}r jeden Build, der sie nicht ausdr{\"u}cklich anfordert: \lstinline[style=fileinline]|CMakeLists.txt|
\item eine Konfigurationsdatei (siehe 6.2.6): \lstinline[style=fileinline]|Kconfig|
\end{itemize}
Die eigentliche Messanwendung ist zudem unterteilt in:

\begin{itemize}
\item eine gemeinsame Basis-Konfiguration (siehe 6.3.1): \lstinline[style=fileinline]|prj.conf|
\item eine Konfigurationsdatei, welche die Auswahl (\lstinline[style=fileinline]|choice BENCH_TYPE|) zwischen den sieben Benchmark-Varianten (die f{\"u}nf Metriken, der Selbsttest und die reine Baseline des Speicherbedarfs) definiert: \lstinline[style=fileinline]|Kconfig|
\item eine CMake-Datei, welche abh{\"a}ngig von dieser Auswahl eine der sechs \lstinline[style=fileinline]|bm_*.c|-Dateien in die zu {\"u}bersetzende Anwendung ein (siehe Kapitel 4.1): \lstinline[style=fileinline]|CMakeLists.txt|
\item einem Einstiegspunkt, welcher die Messbibliothek initialisiert, die \lstinline[style=fileinline]|BENCHMETA|-Kopfzeile ausgibt und anschlie{\ss}end die zur Build-Zeit ausgew{\"a}hlte \lstinline[style=fileinline]|bench_run()|-Implementierung aufruft: \lstinline[style=fileinline]|src/main.c|
\item je eine Implementierung pro Benchmark bzw. Referenzmessung (siehe 4.5): \lstinline[style=fileinline]|src/bm_*.c|
\end{itemize}
\section{Benchmark-Infrastruktur}
Alle sechs ausf{\"u}hrbaren Varianten benutzen die gemeinsame \lstinline[style=fileinline]|bench_run()|-Schnittstelle. Die DWT-basierten Benchmarks teilen Zeitmessung, Statistik und Ausgabe. Der Timer teilt Statistik und Ausgabe, benutzt jedoch eine andere Zeitquelle. Die Heap-Fragmentierung benutzt ein eigenes Zustandsformat. Die Bibliothek vereinheitlicht damit die jeweils gemeinsamen Teile.

\subsection{Warum eine eigene Bibliothek?}
Bevor der Code erkl{\"a}rt wird, wird die Begr{\"u}ndung erl{\"a}utert, warum diese Funktionen nicht einfach in jeder der sechs \lstinline[style=fileinline]|bm_*.c|-Dateien einzeln implementiert wurde. Es gibt daf{\"u}r sowohl einen wissenschaftlichen als auch einen softwaretechnischen Grund.

Der wissenschaftliche Grund ist dabei der wichtigere: Damit die f{\"u}nf Benchmarks und die drei Konfigurationen {\"u}berhaupt miteinander verglichen werden k{\"o}nnen, muss die Art und Weise, wie gemessen wird, f{\"u}r alle exakt identisch sein. W{\"u}rde jede Benchmark-Datei ihre eigene kleine Zeitmess- und Statistiklogik besitzen, k{\"o}nnten kleine Unterschiede in dieser Logik selbst, wie z.B. eine andere Rundung, die eigentlich untersuchten Unterschiede zwischen den Konfigurationen verf{\"a}lschen. Die gemeinsame Bibliothek verringert zudem Unterschiede im Messverfahren. Weitere Einflussgr{\"o}{\ss}en wie z.B. Speicherlayout, Unterbrechungen und Hintergrundaktivit{\"a}t m{\"u}ssen dabei trotzdem beachtet werden.

Der softwaretechnische Grund lautet: Ohne eine gemeinsame Bibliothek m{\"u}sste derselbe code, also Kalibrierung, Statistik, Ausgabeformat, etc. sechsmal fast identisch geschrieben werden. Das kommt mit dem Risiko, dass ein sp{\"a}terer Fehler oder eine Korrektur nur in einer der sechs Kopien nachgezogen wird und die Dateien sich dadurch mit der Zeit voneinander unterscheiden.

\subsection{Die {\"o}ffentliche Schnittstelle}
Die Datei benchlib.h legt fest, was die Bibliothek nach au{\ss}en anbietet, ohne zu zeigen, wie es im Detail funktioniert. Durch diese Kapselung bleiben Einzelheiten der Implementierung, z.B. dass ein Hardware-Z{\"a}hler namens DWT verwendet wird, in den .c-Dateien verborgen. Wer die Bibliothek nutzt, muss nur diese Header Datei verstehen.

In dieser Header Datei wurde das Messprinzip festgelegt. Es wurde ein \glqq{}Sandwich\grqq{}-Muster angewendet. Jeder Benchmark misst nach demselben Schema.

\begin{lstlisting}[style=basecode,language=C]
struct bench_stats {
    uint32_t n;      /* Anzahl Messwerte */
    uint64_t min;    /* Minimum in Zyklen */
    uint64_t max;    /* Maximum in Zyklen */
    double mean;     /* laufender Mittelwert */
    double m2;       /* Summe quadrierter Abweichungen (Welford) */
};

/* Einmalig vor allen Messungen aufrufen: startet den Zyklenzaehler
 *und bestimmt den Messoverhead. */
void bench_timing_init(void);

/* Overhead von zwei direkt aufeinanderfolgender Zaehler-Reads (Zyklen).
 * Vom Messergebnis abziehen. */
uint64_t bench_overhead_cycles(void);

/* Aktuellen Zaehlerstand lesen. */
static inline timing_t bench_now(void)
{
    return timing_counter_get();
}

/* Zyklen zwischen zwei Zaehlerstaenden. */
static inline uint64_t bench_elapsed(timing_t start, timing_t end)
{
    return timing_cycles_get(&start, &end);
}
\end{lstlisting}
Der Z{\"a}hler wird direkt vor und direkt nach diesem Vorgang ausgelesen und bildet anschlie{\ss}end die Differenz. Der zu messende Code befindet sich dabei dazwischen, daher die Bezeichnung \glqq{}Sandwich\grqq{}-Muster.

Bei der Datenstruktur \lstinline[style=fileinline]|struct bench_stats| handelt es sich hierbei um die laufende Statistik f{\"u}r die Benchmarks. Es wurde dabei kein Feld (\emph{array}) implementiert, das alle 1000 Einzelmesswerte speichert, sondern f{\"u}nf zusammenfassende Zahlen:

\begin{itemize}
\item Anzahl der bisherigen Messwerte: \lstinline[style=fileinline]|n|
\item Bisher kleinster gemessener Wert: \lstinline[style=fileinline]|min|
\item Bisher gr{\"o}{\ss}ter gemessener Wert: \lstinline[style=fileinline]|max|
\item Laufender Mittelwert: \lstinline[style=fileinline]|mean|
\item Zwischengr{\"o}{\ss}e f{\"u}r die Streuungsberechnung (siehe 6.2.4): \(M_2\)
\end{itemize}
\begin{lstlisting}[style=basecode,language=C]
/* Laufende Statistik (Welford-Algorithmus: numerisch stabil,
 * braucht keinen Array-Speicher fuer die Einzelwerte). */
struct bench_stats {
	uint32_t n;      /* Anzahl Messwerte */
	uint64_t min;    /* Minimum in Zyklen */
	uint64_t max;    /* Maximum in Zyklen */
	double mean;     /* laufender Mittelwert */
	double m2;       /* Summe quadrierter Abweichungen (Welford) */
};
\end{lstlisting}
Diese Entscheidung wird in 6.2.4 genauer begr{\"u}ndet.

Die Funktionen \lstinline[style=fileinline]|bench_now()| und \lstinline[style=fileinline]|bench_elapsed()| sind im Header selbst als \lstinline[style=fileinline]|static inline| implementiert, nicht in einer .c-Datei. Das hat dabei einen messtechnischen Grund. Ein Funktionsaufruf kostet selbst schon einige CPU-Zyklen, unter anderem f{\"u}r das Anlegen eines Stack-Rahmens und den R{\"u}cksprung. Da hierbei teilweise nur einige Dutzend Zyklen gemessen werden, z.B. die Interrupt Latenz, w{\"u}rde ein normaler Funktionsaufruf f{\"u}r die Zeitmessung selbst schon einen st{\"o}renden Anteil am Messergebnis ausmachen. Daher wurde \lstinline[style=fileinline]|static inline| verwendet. Die Kennzeichnung macht es dem Compiler m{\"o}glich, den Funktionsk{\"o}rper direkt einzubetten. Ob dadurch kein Funktionsaufruf entsteht, muss am erzeugten Maschinencode bzw. Disassembly gepr{\"u}ft werden.

Andere Funktionssignaturen, darunter \lstinline[style=fileinline]|bench_timing_init|, \lstinline[style=fileinline]|bench_overhead_cycles|, \lstinline[style=fileinline]|bench_stats_reset/ add/ stddev|, \lstinline[style=fileinline]|bench_report| und \lstinline[style=fileinline]|bench_run|, werden nur deklariert und nicht direkt implementiert. Sie existieren in den jeweiligen \lstinline[style=fileinline]|.c|-Dateien. Dabei nimmt \lstinline[style=fileinline]|bench_run()| eine Sonderrolle ein. Sie wird hier nur deklariert, aber nirgends in \lstinline[style=fileinline]|benchlib| selbst implementiert. Jede der sechs \lstinline[style=fileinline]|bm_*.c|-Dateien besitzt ihre eigene Version davon und \lstinline[style=fileinline]|main()| ruft sie nach der Initialisierung auf (siehe Kapitel 4.1, Kconfig-choice-Mechanismus).

\subsection{Zeitbasis und Kalibrierung des Messaufwands}
Die Datei \lstinline[style=fileinline]|bench_timing.c| k{\"u}mmert sich um zwei Aufgaben: Die Zeitbasis in Betrieb nehmen und den eigenen Messfehler bestimmen.

\textbf{Zeitbasis starten }Die Funktionen \lstinline[style=fileinline]|timing_init()| und \lstinline[style=fileinline]|timing_start()| sind Funktionen aus Zephyrs eingebauter Timing-Programmierschnittstelle. Auf der hier verwendeten Cortex-M4-Hardware aktivieren sie einen in den Prozessor eingebauten Hardware-Z{\"a}hler namens DWT-Zyklenz{\"a}hler. Dieser z{\"a}hlt bei jedem CPU-Takt eins hoch. Ein Auslesen dieses Z{\"a}hlers vor und nach einem Vorgang und die anschlie{\ss}ende Differenzbildung ergibt dadurch die Anzahl der daf{\"u}r ben{\"o}tigten CPU-Takte. Dabei entspricht ein Takt, bei einer 64 MHz Taktfrequenz, exakt 15,625 Nanosekunden.

\[c_i = c_{i,\mathrm{roh}} - c_{\mathrm{Messaufwand}}\]
\textbf{Kalibrierung des Messaufwands }Selbst das Auslesen des Z{\"a}hlers kostet eine kleine Anzahl an Zyklen. W{\"u}rden diese nicht ber{\"u}cksichtigt werden, dann w{\"u}rde jede Messung um genau diesen kleinen Betrag zu gro{\ss} ausfallen. Um diesen Fehler zu bestimmen, wird der Z{\"a}hler zweimal direkt hintereinander ausgelesen, ohne dass dazwischen irgendein zu messender Vorgang stattfindet. Dabei ist die entstehende Differenz genau der Preis des Messens selbst.

Dieser Vorgang wird 100-mal wiederholt und das Minimum aller 100 Ergebnisse wird als der eigentliche Kalibrierungswert {\"u}bernommen, nicht der Durchschnitt. Das Minimum wird als Sch{\"a}tzung des kleinsten beobachteten Messaufwands benutzt, weil weitere Unterbrechungen einzelne Versuche normalerweise verl{\"a}ngern. Es ist jedoch kein Beweis f{\"u}r einen genau bestimmten wahren Messaufwand. Pipeline-, Prefetch- und Aufrufkontexte k{\"o}nnen sich zwischen Kalibrierung und Benchmark unterscheiden. Die Korrektur verringert dabei den systematischen Leseaufwand, entfernt aber nicht Messpunkt- und Layout-Effekte.

\begin{lstlisting}[style=basecode,language=C]
measured_overhead = UINT64_MAX;
for (int i = 0; i < 100; i++) {
    timing_t t0 = timing_counter_get();
    timing_t t1 = timing_counter_get();
    uint64_t c = timing_cycles_get(&t0, &t1);
    if (c < measured_overhead) {
        measured_overhead = c;
    }
}
\end{lstlisting}
\subsection{Laufende Statistik}
Die Speicherung von 1000 Einzelwerten zu jeweils acht Byte w{\"u}rde 8000 Byte pro gleichzeitig gespeicherter Messreihe brauchen. Da die Benchmarks in getrennten Firmware-Aufbauten ausgef{\"u}hrt werden, summiert sich dieser Bedarf hierbei nicht {\"u}ber alle Konfigurationen hinweg. Die laufende Statistik umgeht trotzdem den zus{\"a}tzlichen Messwertpuffer und h{\"a}lt den daf{\"u}r ben{\"o}tigten Speicher dabei unabh{\"a}ngig von der Anzahl der Wiederholungen. Au{\ss}erdem w{\"u}rde es dabei dem Grundgedanken widersprechen, eine f{\"u}r ressourcenbeschr{\"a}nkte Systeme m{\"o}glichst sparsame Messmethode zu benutzen.

Jedoch entsteht bei dieser Methodik ein Problem. Die Formel f{\"u}r die Varianz, also der Durchschnitt der quadrierten Abweichungen vom Mittelwert, setzt eigentlich voraus, dass der Mittelwert bekannt ist, bevor man die Abweichung berechnen kann. Bei einer laufenden Messung {\"a}ndert sich dabei der Mittelwert mit jedem neuen Wert. Eine andere Methode, durch welche man fortlaufend die Summe aller Werte und die Summe aller quadrierten Werte mitz{\"a}hlt und die Varianz erst am Ende aus beiden Summen berechnet, hat ein numerisches Problem. Bei vielen, {\"a}hnlich gro{\ss}en Werten wird dabei die Differenz von zwei sehr gro{\ss}en, sehr {\"a}hnlichen Zahlen gebildet. Dadurch kann ein Gro{\ss}teil der eigentlich relevanten Nachkommastellen durch Rundungsfehler der Gleitkomma-Arithmetik verloren gehen. Dies wird auch \glqq{}Ausl{\"o}schung\grqq{} genannt. Dies ist ein reales Risiko bei den hier gemessenen Zyklenzahlen.

Die bei dieser Untersuchung verwendete L{\"o}sung ist der Welfords Algorithmus. Der Kerngedanke ist hierbei, den Mittelwert und eine Hilfsgr{\"o}{\ss}e f{\"u}r die Streuung (hier \(M_2\) genannt) bei jedem einzelnen neuen Messwert direkt anzupassen, anstatt sie erst am Ende aus gro{\ss}en Summen zu berechnen. Es wird also berechnet, wie weit der neue Wert vom bisherigen Mittelwert entfernt liegt. Der Mittelwert wird dann um einen kleinen Bruchteil dieser Abweichung in seine Richtung verschoben. Wie gro{\ss} dieser Bruchteil ist, h{\"a}ngt hierbei davon ab, wie viele Werte bereits eingeflossen sind. Bei wenigen Werten, zieht also ein neuer Wert den Mittelwert st{\"a}rker als bei vielen Werten. Anschlie{\ss}end wird ein kleiner Korrekturbetrag zur Streuungsgr{\"o}{\ss}e hinzuaddiert. Das Verfahren umgeht hierbei die problematische Differenz von gro{\ss}en Rohmomentsummen. Rundungsfehler und Grenzen der benutzten Gleitkommadarstellung bleiben dabei dennoch bestehen. Dies passiert zudem bei konstantem, kleinem Speicherbedarf, da nur die f{\"u}nf Felder aus \lstinline[style=fileinline]|struct bench_stats| unabh{\"a}ngig von \(N\) gespeichert werden.

\[n_0 = 0, \qquad \bar{x}_0 = 0, \qquad M_{2,0} = 0\]
\[\delta_n = x_n - \bar{x}_{n-1}\]
\[\bar{x}_n = \bar{x}_{n-1} + \frac{\delta_n}{n}\]
\[M_{2,n} = M_{2,n-1} + \delta_n \left(x_n - \bar{x}_n\right)\]
\[M_{2,n} = M_{2,n-1} + \delta_n \left(x_n - \bar{x}_n\right)\]
Die Benchmarks werden in getrennten Firmware-Images ausgef{\"u}hrt, sodass sich Messwertarrays von den verschiedenen Benchmarks oder Konfigurationen nicht gleichzeitig im RAM summieren w{\"u}rden. Welford umgeht dennoch ein weiteres 8-KiB-Array pro Messfirmware und h{\"a}lt den Speicherbedarf unabh{\"a}ngig von \(N\) konstant. Aus der fortlaufend berechneten Summe quadrierter Abweichungen (\(M_2\)) wird die Stichprobenvarianz durch Division durch \(n-1\) bestimmt. Die anschlie{\ss}end daraus gezogene Quadratwurzel ergibt die Stichprobenstandardabweichung. Die Benutzung von \(n-1\), auch Bessels Korrektur genannt, beachtet hierbei, dass der Mittelwert aus derselben Stichprobe gesch{\"a}tzt wurde. Bei unabh{\"a}ngigen, gleich verteilten Beobachtungen mit endlicher Varianz sch{\"a}tzt die so ausgerechnete Varianz die Populationsvarianz erwartungstreu. Diese Eigenschaft {\"u}bertr{\"a}gt sich jedoch hierbei nicht allgemein auf ihre Quadratwurzel und damit auf die Standardabweichung (Weisstein, o. D., \glqq{}Sample Variance\grqq{}).

\[s^2 = \frac{M_{2,n}}{n-1}\]
\[s = \sqrt{\frac{M_{2,n}}{n-1}}\]
Die Standardabweichung wird deshalb in dieser Arbeit vor allem als deskriptive Kennzahl der Streuung verwendet, die innerhalb eines Messlaufs beobachtet werden.


\subsection{Ergebnisausgabe als CSV}
\begin{lstlisting}[style=basecode,language=C]
#include <zephyr/sys/printk.h>
#include <benchlib.h>

void bench_report(const char *name, const struct bench_stats *s)
{
	/* Alle Werte in CPU-Zyklen; Umrechnung in us erfolgt am PC
	 * (1 Zyklus = 15,625 ns @ 64 MHz, siehe BENCHMETA-Zeile). */
	printk("BENCH,%s,%u,%llu,%llu,%.2f,%.2f\n",
	       name,
	       s->n,
	       (unsigned long long)s->min,
	       (unsigned long long)s->max,
	       s->mean,
	       bench_stats_stddev(s));
}
\end{lstlisting}
Diese Datei gibt am Ende eines Benchmarks genau eine Textzeile pro Metrik aus, in einem einfachen, durch Kommas getrennten Format (CSV, Comma-Separated Values)

BENCH, <Name>, <n>, <min>, \ldots{}

Dieses Format wurde aus zwei Gr{\"u}nden gew{\"a}hlt. Erstens ist es maschinenlesbar. Eine solche Zeile l{\"a}sst sich sp{\"a}ter ohne manuellen Aufwand automatisiert in Tabellenprogramme oder Auswertungsskripte einlesen, ohne dass ein aufw{\"a}ndiger Parser f{\"u}r ein komplexeres Format geschrieben werden m{\"u}sste. Zweitens bleibt es dabei trotzdem menschenlesbar. W{\"a}hrend einer Messsitzung am seriellen Terminal l{\"a}sst sich das Ergebnis direkt am Bildschirm ablesen, ohne ein weiteres Werkzeug zu ben{\"o}tigen. Ein methodisch wichtiger Kommentar direkt im Code zeigt zudem, dass die Ausgabefunktion erst nach Abschluss der Messschleife aufgerufen werden darf. Eine Textausgabe {\"u}ber die serielle Schnittstelle ist selbst ein langsamer Vorgang. W{\"u}rde diese dabei innerhalb der gemessenen 1000 Iterationen aufgerufen werden, dann w{\"u}rde dieser gro{\ss}e, f{\"u}r die Messung irrelevante Zeitaufwand in die Ergebnisse mit einflie{\ss}en. Dadurch w{\"u}rden diese Messungen unbrauchbar gemacht werden. Bei allen ausgegebenen Werten handelt es sich um Rohwerte in CPU-Zyklen. Die Umrechnung in Mikro- bzw. Nanosekunden erfolgt dabei bewusst nicht mehr auf dem Mikrocontroller, sondern erst nachtr{\"a}glich bei der Auswertung am PC. Dabei wird die Taktfrequenz aus der \lstinline[style=fileinline]|BENCHMETA|-Kopfzeile verwendet (siehe Kapitel 4.6.1). Das h{\"a}lt die Firmware selbst einfach und verlagert die eigentliche Interpretation der Daten dorthin, wo sie stattfinden soll.

\subsection{Parametrisierung {\"u}ber Kconfig}
Die Bibliothek wird {\"u}ber ein eigenes Kconfig-Symbol, \lstinline[style=fileinline]|BENCHLIB|, aktiviert, das per \lstinline[style=fileinline]|select TIMING_FUNCTIONS| automatisch die Zephyr-interne Zeitmess-Funktionalit{\"a}t mit einschaltet. Dabei ist dies eine bewusste Entscheidung. Ohne den Zyklenz{\"a}hler w{\"a}re keine der Messfunktionen dieser Bibliothek {\"u}berhaupt nutzbar, weshalb diese Abh{\"a}ngigkeit unbedingt gestaltet wurde, anstatt optional. Zwei weitere Werte sind hierbei konfigurierbar und werden bei jedem Lauf in der \lstinline[style=fileinline]|BENCHMETA|-Kopfzeile mit ausgegeben, damit jede Messreihe aus den Rohdaten heraus nachvollziehbar bleibt:

\lstinline[style=fileinline]|BENCHLIB_ITERATIONS|: Legt die Zahl der ausgewerteten Wiederholungen pro zeitbezogenem Benchmark fest. Der ausgew{\"a}hlte Wert von N=1000 liegt in derselben Gr{\"o}{\ss}enordnung wie der Stichprobenumfang der Leistungs-Auswertung bei Silva et al. (2019, S. 10379).

F{\"u}r unabh{\"a}ngige Beobachtungen mit gleicher Verteilung kann der Standardfehler des Mittelwerts durch die Wurzel aus \(N\) sch{\"a}tzen. Bei \(N=1000\) entspricht dies hierbei ungef{\"a}hr 3,16 Prozent der Stichprobenstandardabweichung. Diese Beziehung ist jedoch keine nachgewiesene Genauigkeitsangabe f{\"u}r die vorliegenden Messreihen. Aufeinanderfolgende Ausf{\"u}hrungen k{\"o}nnen voneinander abh{\"a}ngen, bspw. durch gemeinsamen Systemzustand oder periodische Hintergrundaktivit{\"a}t (Weisstein, o. D., \glqq{}Standard Error\grqq{}).

\lstinline[style=fileinline]|BENCHLIB_WARMUP|: Hierbei handelt es sich um Anzahl der vor der eigentlichen Messung verworfenen \glqq{}Aufw{\"a}rm\grqq{}-Durchl{\"a}ufe. Die allersten Ausf{\"u}hrungen eines Codepfads auf einem Mikrocontroller k{\"o}nnen durch Effekte wie einem Flash-Prefetch bzw. NVMC-Cache, Pipeline, Erstinitialisierung und Thread-Startpfade oder weitere Pipeline-Effekte k{\"u}nstlich langsamer ausfallen als alle folgenden, identischen Durchl{\"a}ufe. Diese ersten zehn Durchl{\"a}ufe werden daher bewusst durchgef{\"u}hrt, aber nicht in die Statistik aufgenommen.

\section{Konfigurationen: Minimal / Logging / (IoT-Stack)}
Alle drei Konfigurationen werden jeweils in einer eigenen \lstinline[style=fileinline]|.conf-Datei| implementiert, die dem Build zusammen mit der gemeinsamen Basis-Konfiguration der Mess-Anwendung {\"u}bergeben wird. Bevor jedoch diese drei Konfigurationsstufen erkl{\"a}rt werden, wird zuerst die gemeinsame Basis erl{\"a}utert, da sie in allen drei Stufen enthalten ist. Dadurch ist sie nicht selbst Teil der untersuchten Variation.

\subsection{Gemeinsame Basis}
\begin{lstlisting}[style=basecode,language=Kconfig]
# --- Zwingend fuer die Messung ---
# Messbibliothek; selektiert TIMING_FUNCTIONS (DWT-Zyklenzaehler)
CONFIG_BENCHLIB=y

# --- Zwingend fuer die Ergebnisausgabe ---
# Synchrone Minimal-Ausgabe (bewusst KEIN Log-Subsystem)
CONFIG_PRINTK=y
# UART-Treiber
CONFIG_SERIAL=y
# Konsolen-Subsystem
CONFIG_CONSOLE=y
# printk -> UART0 (VCOM des J-Link OB)
CONFIG_UART_CONSOLE=y
# %f in printk (mean/stddev); in allen Konfigurationen identisch
# enthalten -> verzerrt den Konfigvergleich nicht
CONFIG_CBPRINTF_FP_SUPPORT=y

# --- Fuer die spaeteren Heap-Benchmarks ---
# 64 KiB System-Heap fuer k_malloc/k_free
CONFIG_HEAP_MEM_POOL_SIZE=65536
\end{lstlisting}
Diese Datei enth{\"a}lt nur das, was f{\"u}r jede Messung technisch zwingend ben{\"o}tigt wird. Dies gilt f{\"u}r jeden Benchmark und jede Konfiguration, die zum Zeitpunkt der Messung aktiv ist. Der gemeinsame Sockel h{\"a}lt hierbei die vorgesehenen Basisfunktionen konstant. Wechselwirkungen mit weiteren aktivierten Subsystemen sind dadurch nicht ausgeschlossen. Der gemeinsame Sockel h{\"a}lt diese Optionen konstant, entfernt aber keine Wechselwirkungen mit weiteren aktivierten Subsystemen oder {\"A}nderungen des Linkerlayouts. Diese Entscheidung ist dabei Teil des konstant gehaltenen Sockels.

\lstinline[style=fileinline]|CONFIG_BENCHLIB=y|

Dies aktiviert das eigenentwickelte Kconfig-Modul der Mess-Bibliothek (siehe 6.2). Dieses eine Flag l{\"o}st intern durch eine Kconfig-select-Abh{\"a}ngigkeit automatisch \lstinline[style=fileinline]|CONFIG_TIMING_FUNCTIONS=|y aus. Dabei handelt es sich um den Zugriff auf Zephyrs Timing-API (Zyklenz{\"a}hler der Hardware, siehe 6.2). Diese Flag erlaubt es, die Zeitmessfunktionen zu benutzen, welche in den Benchmarks verwendet werden.

\lstinline[style=fileinline]|CONFIG_PRINTK=y, CONFIG_SERIAL=y, CONFIG_CONSOLE=y, CONFIG_UART_CONSOLE=y|

Diese Optionen stellen die vorgesehene Konsolenausgabe der Messanwendung bereit. {\"U}ber den virtuellen COM-Port vom Onboard-Debugger k{\"o}nnen die ausgegebenen Messdaten am Host-Rechner direkt empfangen werden. Die Benutzung derselben Ausgabefunktion und derselben Schnittstelle legt jedoch nicht alle internen Verarbeitungsschritte fest. Vor allem kann ein aktiviertes Logging-Subsystem \lstinline[style=fileinline]|printk()|-Nachrichten {\"u}bernehmen. Ob dies im jeweiligen Messlauf passiert, ist anhand der zusammengef{\"u}hrten Build-Konfiguration zu pr{\"u}fen. In der Logging-Stufe stellt \lstinline[style=fileinline]|CONFIG_LOG_PRINTK=n| sicher, dass \lstinline[style=fileinline]|printk()| nicht zum Logging-Frontend umgeleitet wird. (Zephyr Project, 2024, \glqq{}Logging\grqq{}, Abschnitt \glqq{}Printk\grqq{}).

\lstinline[style=fileinline]|CONFIG_CBPRINTF_FP_SUPPORT=y|

Zephyrs \lstinline[style=fileinline]|printk()| unterst{\"u}tzt in der Grundeinstellung aus Ressourcengr{\"u}nden keine Gleitkommazahlen (\lstinline[style=fileinline]|%f|), da die daf{\"u}r ben{\"o}tigte Software-Gleitkomma-Formatierung Flash-Speicher kostet. Da die Mess-Statistik aber Mittelwert und Standardabweichung als Kommazahlen ausgibt, muss dieses Flag gesetzt sein. Dabei ist es wichtig zu erw{\"a}hnen, dass dies in der Messmethodik in allen drei Konfigurationsstufen implementiert ist. Dadurch ist der verursachte Flash-Speicher Verbrauch in jeder Messung gleich hoch und verf{\"a}lscht dadurch nicht den Vergleich. Dieser Verbrauch ist also Teil des konstanten Sockels und nicht der untersuchten Variablen.

\lstinline[style=fileinline]|CONFIG_HEAP_MEM_POOL_SIZE=65536|

Dies reserviert in jedem Benchmark-Profil einen 64-KiB-Systemheap. Dieser wird direkt vom Allokationsbenchmark benutzt, steht an sich aber auch anderen aktivierten Subsystemen zur Verf{\"u}gung. Die Reservierung beeinflusst den statischen RAM-Report und wird deshalb als Teil des gemeinsamen Sockels ausgewiesen. Der Fragmentierungsbenchmark benutzt davon unabh{\"a}ngig einen eigenen 8-KiB-\lstinline[style=fileinline]|k_heap|.

\subsection{Minimal-Konfiguration}
\begin{lstlisting}[style=basecode,language=Kconfig]
CONFIG_BOOT_BANNER=n
CONFIG_TIMESLICING=n
CONFIG_ASSERT=n
CONFIG_LOG=n
CONFIG_THREAD_NAME=n
CONFIG_DEBUG_THREAD_INFO=n
CONFIG_ARM_MPU=n
CONFIG_HW_STACK_PROTECTION=n
CONFIG_USE_SEGGER_RTT=n
CONFIG_GPIO=n
\end{lstlisting}
Diese Konfiguration ist die erste und simpelste Untersuchungsstufe. Ziel dieser Konfiguration ist es, einen RTOS-Kern zu erhalten, der nur das enth{\"a}lt, was f{\"u}r die Mess-Anwendung unbedingt ben{\"o}tigt wird. Jedes Merkmal, das zus{\"a}tzlich von Zephyr oder dem verwendeten Board standardm{\"a}{\ss}ig aktiviert w{\"u}rde, wird hierbei bewusst deaktiviert. Dadurch kann die Minimalkonfiguration als Referenzstufe (Baseline) fungieren, um sp{\"a}ter die Logging- und IoT-Stufe mit ihr vergleichen zu k{\"o}nnen. Dabei ist zu beachten, dass nicht jedes Flag in dieser Datei einen von Zephyr aktivierten Standardwert deaktiviert. Ein Teil der Flags ist bereits die Zephyr-Voreinstellung und wird hierbei explizit dokumentiert, damit es nachvollziehbar ist, dass diese Entscheidungen bewusst getroffen sind. Andere Flags wiederum {\"u}berschreiben einen abweichenden Standard- bzw. Board-Wert.

\lstinline[style=fileinline]|CONFIG_BOOT_BANNER=n|

Zephyr gibt beim Start standardm{\"a}{\ss}ig eine Textzeile mit Versionsinformationen {\"u}ber die Konsole aus, den sogenannten \emph{Boot-Banner}. Diese Ausgabe kostet zus{\"a}tzlichen Flash-Speicher sowohl f{\"u}r die Zeichenkette als auch f{\"u}r die Laufzeit beim Start. F{\"u}r die Messungen ist diese Ausgabe unn{\"o}tig. Durch das Deaktivieren entsteht dabei ein praktischer Nebeneffekt. Falls der Boot Banner trotz Deaktivierung ausgegeben wird, zeigt dies, dass dieses Konfigurationsfragment beim Build nicht angewendet wurde.

\lstinline[style=fileinline]|CONFIG_TIMESLICING=n|

Zephyrs Scheduler kann so konfiguriert werden, dass er gleichprioren Threads per Zeitscheiben-Rotation abwechselnd Rechenzeit zuteilt (Round-Robin), indem er einen laufenden Thread nach Ablauf eines Zeitfensters unterbricht, auch wenn dieser noch nicht freiwillig abgeben wollte. F{\"u}r die Benchmarks ist dieses Verhalten jedoch unerw{\"u}nscht, da eine solche unvorhersehbare Unterbrechung einen St{\"o}rfaktor in der Zeitmessung darstellt, der nichts mit der untersuchten Metrik zu tun hat. Dadurch w{\"u}rden die einzelnen Messwerte durch sogenannte\emph{ Ausrei{\ss}er} nach oben verf{\"a}lscht werden. Diese Entscheidung folgt dabei derselben Vorgehensweise wie Zephyrs eigener Referenz-Benchmark (\lstinline[style=fileinline]|latency_measure|), der ebenfalls ohne Zeitscheiben-Rotation misst.

\lstinline[style=fileinline]|CONFIG_ASSERT=n|

Dies aktiviert bzw. deaktiviert die im Kernel eingebaute Laufzeitpr{\"u}fung (Assertions), die im Falle eines Fehlers das System kontrolliert anhalten. Dieser Wert entspricht dabei bereits der Zephyr-Standardeinstellung f{\"u}r den hier verwendeten Build-Typ. Dies wird hierbei explizit gesetzt, um die bewusst getroffene Entscheidung zu dokumentieren, anstatt sich auf einen unsichtbaren Standardwert zu verlassen.

\lstinline[style=fileinline]|CONFIG_LOG=n|

Dies deaktiviert das Zephyr-Logging-Subsystem vollst{\"a}ndig. Dieses Flag ist die entscheidende Kontrollvariable dieser Untersuchung. Es handelt sich hierbei um den einzigen inhaltlichen Unterschied zur Logging-Konfiguration (siehe 6.3.3). Die beobachteten Unterschiede werden hierbei zwischen zwei Profilen untersucht, die sich durch die Aktivierung des Logging-Subsystems unterscheiden. Ihre technische Ursache muss jedoch noch anhand von der aktiven Konfiguration und des erzeugten Programms beurteilt werden.

\lstinline[style=fileinline]|CONFIG_THREAD_NAME=n|

Dies deaktiviert die Speicherung eines lesbaren Namensstrings pro Thread (z.B. f{\"u}r Debugging-Zwecke). Jeder gespeicherte Name kostet extra RAM-Platz pro angelegten Thread. Diese Flag ist deaktiviert, da dies f{\"u}r die Messungen nicht ben{\"o}tigt wird.

\lstinline[style=fileinline]|CONFIG_DEBUG_THREAD_INFO=n|

Dies deaktiviert zus{\"a}tzliche, f{\"u}r Debugger bestimmte Metadaten {\"u}ber die im System laufenden Threads (z.B. f{\"u}r die Anzeige im Eclipse-/ GDB-Thread-Debugging). Auch dieses Merkmal dient der Entwicklungsunterst{\"u}tzung und wird f{\"u}r das Ausf{\"u}hren der Mess-Anwendung nicht ben{\"o}tigt.

\lstinline[style=fileinline]|CONFIG_ARM_MPU=n & CONFIG_HW_STACK_PROTECTION=n|

Diese Flags deaktivieren den MPU-basierten Schutz in allen drei Benchmark-Profilen. Bei aktivem Stack-Guard k{\"o}nnen auch Supervisor-Threads MPU-Regionen brauchen. Userspace-Memory-Domains sorgen f{\"u}r weitere MPU-Konfiguration. Da keine dieser Varianten gemessen wurde, l{\"a}sst sich der vermiedene Kontextwechselaufwand nicht quantifizieren. Die Messergebnisse gelten daher f{\"u}r die ausgew{\"a}hlte Konfiguration ohne MPU-basierten Stackschutz.

\lstinline[style=fileinline]|CONFIG_USE_SEGGER_RTT=n|

Dieses Flag wurde durch eine Untersuchung w{\"a}hrend der Einrichtung gesetzt. Dabei wurde ein Speicherbericht (\lstinline[style=fileinline]|ram_report|) erstellt, welcher zeigte, dass der Treiber f{\"u}r das SEGGER-RTT-Protokoll ungefragt in das Firmware-Abbild eingebunden wurde, obwohl die Ausgabe der Mess-Anwendung nur {\"u}ber UART erfolgt. Bei diesem Protokoll handelt es sich um einen alternativen, im Debugger integrierten Ausgabekanal. Dies liegt daran, dass es sich um einen Board-Standardwert des verwendeten nRF52840-DK-Boards handelt, keinen allgemeinen Zephyr-Standardwert. Er wurde daher explizit deaktiviert, nachdem sein ungenutzter Speicherverbrauch (ca. 2,2 KiB Flash und RAM f{\"u}r einen internen Puffer) im Speicherbericht angezeigt wurde.

\lstinline[style=fileinline]|CONFIG_GPIO=n|

Dieses Flag wurde ebenfalls gesetzt, nachdem ein Board-Standardwert entdeckt wurde. Der GPIO-Treiber (allgemeine digitale Ein-/ Ausgabepins) wurde vom Board-Definitionsdatensatz standardm{\"a}{\ss}ig aktiviert, obwohl die Mess-Anwendung keine GPIO-Funktionalit{\"a}t verwendet. Es wird nur die serielle Schnittstelle, die {\"u}ber eine eigene, unabh{\"a}ngige Pin-Konfiguration verf{\"u}gt, verwendet. Auch dies wurde durch derselben \lstinline[style=fileinline]|ram_report|-Pr{\"u}fung herausgefunden.

\subsection{Logging-Konfiguration}
\begin{lstlisting}[style=basecode,language=Kconfig]
CONFIG_LOG=y
CONFIG_LOG_PRINTK=n
CONFIG_LOG_MODE_DEFERRED=y
CONFIG_LOG_BUFFER_SIZE=1024
CONFIG_LOG_DEFAULT_LEVEL=3
CONFIG_LOG_BACKEND_UART=y
CONFIG_LOG_FUNC_NAME_PREFIX_DBG=n
\end{lstlisting}
Diese Konfiguration entspricht exakt der Minimal-Konfiguration mit genau einer gro{\ss}en {\"A}nderung: \lstinline[style=fileinline]|CONFIG_LOG| wird von \lstinline[style=fileinline]|n| auf \lstinline[style=fileinline]|y| gesetzt, wodurch das Zephyr-Logging-Subsystem aktiviert wird. Wichtig f{\"u}r die Reproduzierbarkeit ist hierbei, dass die Datei \lstinline[style=fileinline]|logging.conf| nicht die Einstellungen aus \lstinline[style=fileinline]|minimal.conf| als Text enth{\"a}lt, sondern nur das inhaltliche Delta ist. Damit diese Konfigurationsstufe aus der Minimal- und Logging-Konfiguration besteht, m{\"u}ssen beim Bauen beide Dateien gemeinsam {\"u}bergeben werden. Dies erfolgt durch den Befehl:

\lstinline[style=fileinline]|-DEXTRA_CONF_FILE="minimal.conf;logging.conf"| (siehe Kapitel 4.7)

Nur so bleiben auch in dieser Konfigurationsstufe die, in der Minimal-Konfiguration, deaktivierten Merkmale weiterhin deaktiviert. Der einzige inhaltliche Unterschied wird somit auf das Logging-Merkmal beschr{\"a}nkt.

\lstinline[style=fileinline]|CONFIG_LOG=y|

Aktiviert das Logging-Subsystem als Ganzes. Dies ist die eigentliche unabh{\"a}ngige Variable dieser Konfigurationsstufe.

\lstinline[style=fileinline]|CONFIG_LOG_MODE_DEFERRED=y|

Zephyrs Logging-Subsystem kennt hierbei zwei Betriebsarten. Im \emph{Immediate}-Modus w{\"u}rde jeder einzelne Log-Aufruf sofort oder synchron die komplette Textformatierung durchf{\"u}hren und {\"u}ber die serielle Schnittstelle ausgeben, bevor der aufrufende Code fortgesetzt wird. F{\"u}r die Untersuchung wurde jedoch die zweite Betriebsart, der \emph{Deferred}-Modus, verwendet. Hierbei schreibt ein Log-Aufruf stattdessen nur die Rohdaten (Nachricht und Argumente) in einen Ringpuffer. Ein separater, niedriger priorisierter Log-Thread entnimmt dann diese Eintr{\"a}ge asynchron im Hintergrund. Diese werden dann erst dort in Text umformatiert und anschlie{\ss}end ausgegeben. Diese Entscheidung wurde bewusst getroffen. Der Immediate-Modus w{\"u}rde direkt im gemessenen Codepfad selbst unvorhersehbar lange, synchrone UART-Schreibvorg{\"a}nge ausl{\"o}sen und dadurch die Zeitmetriken, die untersucht werden, verf{\"a}lschen. Der Deferred-Modus entspricht dem in der Praxis {\"u}blichen Einsatzszenario von Logging in eingebetteten Systemen und ist Zephyrs Standardeinstellung.

\lstinline[style=fileinline]|CONFIG_LOG_BUFFER_SIZE=1024|

Dies legt die Gr{\"o}{\ss}e des eben genannten Ringpuffers auf 1024 Byte fest. Dabei handelt es sich um einen Zephyr-Standardwert, welcher {\"u}bernommen wurde. Dieser Speicherbereich existiert nur in der Logging-Konfiguration und ist damit ein direkt messbarer RAM-Mehrverbrauch, welcher sich der Logging-Funktion zuordnen l{\"a}sst.

\lstinline[style=fileinline]|CONFIG_LOG_DEFAULT_LEVEL=3|

Dies legt f{\"u}r Module ohne eigene Festlegung das Informationsniveau als Standardfilter fest. Daraus ergibt sich dann jedoch keine bestimmte Anzahl von Log-Nachrichten. Die echte Logging-Last h{\"a}ngt zudem davon ab, welche Log-Aufrufe ausgef{\"u}hrt werden, wie h{\"a}ufig sie auftreten und welche Inhalte sie verarbeiten.  (Zephyr Project, 2024, \glqq{}Logging\grqq{}, Abschnitt \glqq{}Global Kconfig Options\grqq{}).

\lstinline[style=fileinline]|CONFIG_LOG_BACKEND_UART=y|

Dies legt fest, {\"u}ber welchen physischen Ausgabekanal der Log-Thread seine formatierten Nachrichten ausgibt. Hierbei wird derselbe UART-Kanal verwendet, {\"u}ber den auch die printk()-Ergebniszeilen der Benchmarks laufen. Diese Wahl stellt sicher, dass keine, in der Minimal-Konfiguration nicht vorhandene Hardware-Schnittstelle (z.B SEGGER-RTT) f{\"u}r das Logging aktiviert wird.

\lstinline[style=fileinline]|CONFIG_LOG_FUNC_NAME_PREFIX_DBG=n|

Dieses Flag deaktiviert die automatische Voranstellung des Funktionsnamens vor jede Debug-Log-Nachricht. Dieses Merkmal betrifft nur die Textformatierung der Ausgabe und keine der Benchmark-relevanten Zeitmetriken. Es wurde daher deaktiviert, um den daraus entstehenden Flash-Speicherverbrauch zu vermeiden.

\subsection{IoT-Stack-Konfiguration}
\begin{lstlisting}[style=basecode,language=Kconfig]
CONFIG_NETWORKING=y
CONFIG_NET_SOCKETS=y
CONFIG_NET_SOCKETS_POSIX_NAMES=y
CONFIG_NET_IPV6=y
CONFIG_NET_UDP=y
CONFIG_NET_TCP=n
CONFIG_NET_L2_OPENTHREAD=y
CONFIG_OPENTHREAD_MTD=y
CONFIG_OPENTHREAD_MANUAL_START=y
CONFIG_COAP=y
CONFIG_NET_SOCKETS_SOCKOPT_TLS=y
CONFIG_NET_SOCKETS_ENABLE_DTLS=y
CONFIG_MBEDTLS=y
CONFIG_MBEDTLS_TLS_VERSION_1_2=y
CONFIG_MBEDTLS_ENABLE_HEAP=y
CONFIG_MBEDTLS_HEAP_SIZE=16384
CONFIG_MBEDTLS_KEY_EXCHANGE_PSK_ENABLED=y
CONFIG_MBEDTLS_PSK_MAX_LEN=32
\end{lstlisting}
Die Benchmark-Stufe kombiniert \lstinline[style=fileinline]|minimal.conf| mit \lstinline[style=fileinline]|iot_stack.conf| und \lstinline[style=fileinline]|iot_stack.overlay|. Aktiviert werden Networking, IPv6, UDP, OpenThread MTD, CoAP und DTLS/ PSK {\"u}ber mbedTLS.

\lstinline[style=fileinline]|CONFIG_NETWORKING=y|

Diese Option aktiviert Zephyrs allgemeine Netzwerkunterst{\"u}tzung. Sie bildet die Grundlage f{\"u}r Link-Layer, Netzwerkpuffer und h{\"o}here Protokolle wie IPv6. Die Option aktiviert allein noch keine bestimmte Netzwerkkommunikation, stellt aber die gemeinsamen Strukturen bereit, auf denen die nachfolgenden Netzwerkoptionen aufbauen.

\lstinline[style=fileinline]|CONFIG_NET_SOCKETS=y|

Diese Option aktiviert Zephyrs BSD-Socket-Schnittstelle. Anwendungen k{\"o}nnen dadurch Netzwerkkommunikation {\"u}ber Funktionen wie \lstinline[style=fileinline]|socket()|, \lstinline[style=fileinline]|bind()|, \lstinline[style=fileinline]|recvfrom()| und \lstinline[style=fileinline]|sendto()| umsetzen. Die Socket-Architektur wird bei der Implementierung des CoAP-Servers in Kapitel 6.4.2 benutzt. Die allgemeine Rolle der Socket-Schnittstelle wird deshalb hier nicht nochmal erl{\"a}utert (Zephyr Project, 2024, \glqq{}BSD Sockets\grqq{}, Abschnitt \glqq{}Overview\grqq{}).

\lstinline[style=fileinline]|CONFIG_NET_SOCKETS_POSIX_NAMES=y|

Zephyr versieht seine Socket-Funktionen standardm{\"a}{\ss}ig mit dem Pr{\"a}fix \lstinline[style=fileinline]|zsock_|, um Namenskonflikte zu umgehen. Diese Option stellt zudem die bekannten POSIX-Namen wie \lstinline[style=fileinline]|socket()|, \lstinline[style=fileinline]|bind()|, \lstinline[style=fileinline]|recvfrom()| und \lstinline[style=fileinline]|sendto()| bereit, ohne daf{\"u}r die komplette POSIX-Unterst{\"u}tzung zu aktivieren.

Die Option ist in Zephyr v3.7.2 schon als veraltet gekennzeichnet und wird aufgrund der regul{\"a}ren POSIX-Socket-API nur noch aus Kompatibilit{\"a}tsgr{\"u}nden angeboten. Sie wird hier benutzt, weil der vorhandene Anwendungscode diese Funktionsnamen benutzt (Zephyr Project, 2024, \glqq{}BSD Sockets\grqq{}; Zephyr Project, o. D., Kconfig NET\_SOCKETS\_POSIX\_NAMES, Tag v3.7.2).

\lstinline[style=fileinline]|CONFIG_NET_IPV6=y|

Diese Option aktiviert den IPv6-Teil vom Zephyr-Netzwerkstack. IPv6 ist f{\"u}r das benutzte Thread-Netz erforderlich, da Thread als IPv6-basiertes Mesh-Protokoll arbeitet. Die Verbindung zwischen Thread, IPv6, 6LoWPAN und IEEE 802.15.4 wurde bereits in Kapitel 2.7.3 beschrieben.

\lstinline[style=fileinline]|CONFIG_NET_UDP=y|

Diese Option aktiviert UDP als verbindungsloses Transportprotokoll. CoAP benutzt UDP. Auch die DTLS-Absicherung arbeitet in dieser Konfiguration auf Datagrammen. Die Protokollbeziehung zwischen CoAP, UDP und DTLS wurde bereits in Kapitel 2.7.1 und 2.7.2 erl{\"a}utert.

\lstinline[style=fileinline]|CONFIG_NET_TCP=n|

TCP wird deaktiviert, da die Benchmark-Konfiguration keine TCP-Anwendung und kein TLS {\"u}ber einen TCP-Datenstrom benutzt. CoAP und DTLS arbeiten hier auf UDP. Die Deaktivierung verhindert, dass nicht gebrauchte TCP-Funktionen in das Firmware-Abbild aufgenommen werden.

\lstinline[style=fileinline]|CONFIG_NET_L2_OPENTHREAD=y|

Diese Option bindet OpenThread als Layer-2-Implementierung in Zephyrs Netzwerkstack ein. Zephyr verbindet dabei den externen OpenThread-Stack mit seiner IEEE-802.15.4-Treiberschnittstelle. Die allgemeinen Eigenschaften des Thread-Protokolls wurden bereits in Kapitel 2.7.3 erkl{\"a}rt.

Die Aktivierung zieht weitere Abh{\"a}ngigkeiten nach sich. Dazu geh{\"o}ren unter anderem die IEEE-802.15.4-Unterst{\"u}tzung, OpenThread selbst, C++-Laufzeitanteile und eine Entropiequelle. Der Speicherbedarf der IoT-Stufe ist daher nicht nur den 18 sichtbaren Zeilen des Konfigurationsfragments zuzuordnen, sondern beinhaltet auch die von Kconfig aufgel{\"o}sten Abh{\"a}ngigkeiten (Zephyr Project, 2023, \glqq{}Thread protocol\grqq{}; Zephyr Project, o. D., Kconfig NET\_L2\_OPENTHREAD, Tag v3.7.2).

\lstinline[style=fileinline]|CONFIG_OPENTHREAD_MTD=y|

Diese Option w{\"a}hlt die Ger{\"a}tekategorie Minimal Thread Device (MTD). Ein MTD beteiligt sich nicht als Router an der Weiterleitung fremder Pakete. Dadurch ist der Funktionsumfang niedriger als bei einem Full Thread Device. Diese Auswahl entspricht dem ressourcenbeschr{\"a}nkten Endger{\"a}t, das durch die IoT-Sensoranwendung dargestellt werden soll.

\lstinline[style=fileinline]|CONFIG_OPENTHREAD_MANUAL_START=y|

Diese Option verhindert, dass OpenThread w{\"a}hrend der Initialisierung automatisch versucht, einem vorhandenen Thread-Netzwerk beizutreten. Das Stack m{\"u}sste stattdessen {\"u}ber seine API oder eine passende Kommandozeilenschnittstelle gestartet werden.

F{\"u}r die Benchmark-Anwendung erfolgt ein solcher Start nicht. Dadurch werden keine fortlaufenden Beitrittsversuche ohne vorhandenen Thread-Koordinator erzeugt. Die Messung beinhaltet somit die eingebundene OpenThread-Software und ihre statischen Ressourcen, aber keinen aktiven Netzwerkbeitritt oder Mesh-Verkehr.

Dieser Punkt unterscheidet die Benchmark-Konfiguration von \lstinline[style=fileinline]|iot_sensor|. Dort ist \lstinline[style=fileinline]|CONFIG_OPENTHREAD_MANUAL_START=n| gesetzt, weil die Demonstrationsanwendung am Netzwerkbetrieb teilnehmen soll (Zephyr Project, o. D., Kconfig OPENTHREAD\_MANUAL\_START, Tag v3.7.2).

\lstinline[style=fileinline]|CONFIG_COAP=y|

Diese Option bindet Zephyrs CoAP-Bibliothek ein. Sie stellt Funktionen zum Erzeugen und Parsen von CoAP-Paketen und zur Verarbeitung von Ressourcen und Optionen bereit. Die Bibliothek erstellt selbst keine Sockets. Die Anwendung stellt den Kommunikationskanal {\"u}ber die Socket-API bereit und {\"u}bergibt dabei die empfangenen Puffer zur Verarbeitung an die CoAP-Bibliothek.

Das Protokoll selbst wurde bereits in Kapitel 2.7.1 erl{\"a}utert. Die bestimmte Verwendung f{\"u}r die Ressource \lstinline[style=fileinline]|/sensor| wird in Kapitel 6.4.2 beschrieben (Zephyr Project, 2023, \glqq{}CoAP\grqq{}, Abschnitt \glqq{}Overview\grqq{}).

In der Benchmark-Anwendung wird keine CoAP-Ressource aufgerufen. Die Option dient dort dazu, die Bibliothek als Teil des IoT-Softwareprofils in den Build einzubeziehen.

\lstinline[style=fileinline]|CONFIG_NET_SOCKETS_SOCKOPT_TLS=y|

Diese Option erweitert die Socket-Schnittstelle um TLS-/ DTLS-spezifische Protokolltypen und Socket-Optionen. Dazu geh{\"o}ren bspw. \lstinline[style=fileinline]|TLS_SEC_TAG_LIST| zur Auswahl von registrierten Zugangsdaten und \lstinline[style=fileinline]|TLS_DTLS_ROLE| zur Festlegung der Client- oder Serverrolle.

Die allgemeine DTLS-/ PSK-Funktionsweise wurde schon in Kapitel 2.7.2 erl{\"a}utert. Die bestimmte Registrierung eines PSK und seiner Identit{\"a}t {\"u}ber \lstinline[style=fileinline]|tls_credential_add()| und deren Zuordnung zum Socket wird in Kapitel 6.4.2 beschrieben (Zephyr Project, 2024, \glqq{}BSD Sockets\grqq{}, Abschnitt \glqq{}Secure Sockets\grqq{}).

\lstinline[style=fileinline]|CONFIG_NET_SOCKETS_ENABLE_DTLS=y|

Diese Option erweitert die sichere Socket-Schnittstelle um DTLS-Unterst{\"u}tzung. Ohne sie w{\"a}re {\"u}ber \lstinline[style=fileinline]|CONFIG_NET_SOCKETS_SOCKOPT_TLS| standardm{\"a}{\ss}ig nur TLS auf einem TCP-Datenstrom verf{\"u}gbar. Die Option w{\"a}hlt dabei f{\"u}r Zephyrs nativen Netzwerkstack gleichzeitig die gebrauchte DTLS-Unterst{\"u}tzung in mbedTLS aus.

Obwohl die Funktion in das Benchmark-Firmware-Abbild eingebunden ist, erzeugt \lstinline[style=fileinline]|bench| keinen DTLS-Socket und f{\"u}hrt keinen Handshake aus. Gemessen wird daher nicht die Laufzeit einer DTLS-Verbindung, sondern der Einfluss des eingebundenen Softwareprofils.

\lstinline[style=fileinline]|CONFIG_MBEDTLS=y|

Diese Option bindet die mbedTLS-Kryptografiebibliothek ein. Zephyrs native TLS-/ DTLS-Sockets benutzen deren kryptografische und protokollspezifische Funktionen. mbedTLS wird au{\ss}erdem von OpenThread f{\"u}r verschiedene Sicherheitsfunktionen benutzt.

\lstinline[style=fileinline]|CONFIG_MBEDTLS_TLS_VERSION_1_2=y|

Diese Option aktiviert die Protokollunterst{\"u}tzung f{\"u}r TLS 1.2 in mbedTLS. In der hier benutzten Zephyr-Version geh{\"o}rt diese Einstellung zu den Voraussetzungen f{\"u}r die aktivierte DTLS-Unterst{\"u}tzung. Die Benennung bezieht sich dabei auf die gemeinsame TLS-1.2-Protokollbasis. Im Socket wird dann anschlie{\ss}end DTLS 1.2 verwendet.

\lstinline[style=fileinline]|CONFIG_MBEDTLS_ENABLE_HEAP=y|

Diese Option erlaubt mbedTLS die Verwendung eines eigenen globalen Speicherbereichs f{\"u}r dynamische Anforderungen. Der mbedTLS-Heap soll dabei von dem {\"u}ber \lstinline[style=fileinline]|CONFIG_HEAP_MEM_POOL_SIZE| bereitgestellten Zephyr-Systemheap unterschieden werden. Zephyr initialisiert diesen Speicherbereich beim Systemstart automatisch, solange die regul{\"a}re mbedTLS-Initialisierung aktiv ist (Zephyr Project, o. D., Kconfig MBEDTLS\_ENABLE\_HEAP, Tag v3.7.2).

\lstinline[style=fileinline]|CONFIG_MBEDTLS_HEAP_SIZE=16384|

Diese Option legt die Gr{\"o}{\ss}e des mbedTLS-Heaps auf 16.384 Byte bzw. 16 KiB fest. Dieser reservierte Speicherbereich ist ein Teil des RAM-Bedarfs der IoT-Konfiguration.

Die Gr{\"o}{\ss}e ist eine bewusste Konfigurationsentscheidung f{\"u}r diese Untersuchung. Sie ist kein allgemeiner mbedTLS-Standardwert. Der gebrauchte Speicher h{\"a}ngt von den benutzten Protokollen, Puffergr{\"o}{\ss}en, Cipher Suites und gleichzeitig aktiven Verbindungen ab (Zephyr Project, o. D., Kconfig MBEDTLS\_HEAP\_SIZE, Tag v3.7.2).

\lstinline[style=fileinline]|CONFIG_MBEDTLS_KEY_EXCHANGE_PSK_ENABLED=y|

Diese Option aktiviert PSK-basierte Schl{\"u}sselaustauschverfahren in mbedTLS. Dadurch kann die in Kapitel 2.7.2 beschriebene Absicherung mit einem vorher ausgesuchten gemeinsamen Schl{\"u}ssel benutzt werden.

Die Option stellt nur die gebrauchte kryptografische Funktionalit{\"a}t bereit. In \lstinline[style=fileinline]|bench| werden keine Zugangsdaten registriert und kein Handshake ausgef{\"u}hrt. Die Verwendung eines PSK erfolgt erst in \lstinline[style=fileinline]|iot_sensor|, wie in Kapitel 6.4.2 beschrieben.

\lstinline[style=fileinline]|CONFIG_MBEDTLS_PSK_MAX_LEN=32|

Diese Option begrenzt die maximale L{\"a}nge eines Pre-Shared Keys auf 32 Byte. Sie definiert damit eine Kapazit{\"a}tsgrenze f{\"u}r PSK-basierte Verbindungen, erzeugt aber selbst keinen Schl{\"u}ssel.

Der in der Demonstrationsanwendung benutzte PSK muss innerhalb dieser Grenze liegen. \lstinline[style=fileinline]|MBEDTLS_PSK_MAX_LEN| ist der Symbolname f{\"u}r die benutzte Zephyr-/ mbedTLS-Version.

\textbf{Devicetree-Overlay f{\"u}r die Entropiequelle}

Zum \lstinline[style=fileinline]|iot_stack.conf| wird zudem \lstinline[style=fileinline]|iot_stack.overlay| eingebunden:

\begin{lstlisting}[style=basecode,language=DTS]
/ {
    chosen {
        zephyr,entropy = &rng;
    };
};
\end{lstlisting}
Das Overlay weist Zephyr den Hardware-Zufallszahlengenerator rng als Entropiequelle zu. OpenThread und die kryptografischen Funktionen brauchen dabei Zufallswerte, bspw. f{\"u}r Protokoll- und Sicherheitsoperationen. Die Entropiequelle wird daher {\"u}ber Devicetree festgelegt, weil sie eine Zuordnung zu einem bestimmten Hardwareger{\"a}t beschreibt. Sie ist deshalb nicht als weitere \lstinline[style=fileinline]|CONFIG_|-Option in \lstinline[style=fileinline]|iot_stack.conf| enthalten.

Das Overlay wird nur f{\"u}r die IoT-Stack-Konfiguration {\"u}bergeben. Minimal- und Logging-Konfiguration werde dadurch nicht von dieser Hardwarezuordnung beeinflusst.

\section{Realistische Anwendung}
F{\"u}r diese Arbeit wurde der Quellcode von einer praxisnahen Demonstrationsanwendung entwickelt. Nach einem erfolgreichem Build, Flash- und Ende-zu-Ende-Test soll sie einen {\"u}ber DTLS gesch{\"u}tzten CoAP-Sensordienst und einen Firmware-Update-Pfad demonstrieren.

\subsection{Zweck und Sensor-Simulation}
F{\"u}r diese Arbeit wurde kein externer Sensor an das nRF52840 DK angeschlossen, um die Reproduzierbarkeit nicht von weiterer, nicht mitgelieferter Hardware zu beeinflussen. Daher wird der interne Die-Temperatursensor des nRF52840 DK ausgelesen. Dieser liefert die aktuelle Chiptemperatur, die dann anschlie{\ss}end {\"u}ber den CoAP-Server zur Verf{\"u}gung gestellt wird. Die f{\"u}r OpenThread gebrauchte Entropiequelle wird dabei {\"u}ber ein board-spezifisches Devicetree-Overlay umgestellt. Dies erfolgt nach demselben Prinzip wie bei der in 6.3.4 beschriebenen IoT-Stack-Konfigurationsstufe.

\subsection{CoAP-Server {\"u}ber DTLS/ PSK}
Die Datei \lstinline[style=fileinline]|coap_server.c| implementiert hierbei einen CoAP-Server, der {\"u}ber Zephyrs BSD-Socket-Schnittstelle und die in Zephyr enthaltene CoAP-Paketbibliothek genau eine Ressource zur Verf{\"u}gung stellt. Dabei handelt es sich um eine GET-Anfrage auf dem Pfad \lstinline[style=fileinline]|/sensor|, welche den aktuellen Temperaturwert als Textnutzlast zur{\"u}ckliefert. Die Absicherung des Servers folgt dabei genau dem in Abschnitt 2.7.2 erw{\"a}hntem Verfahren. Anstatt dass ein UDP-Socket benutzt wird, wird beim Anlegen des Sockets das Protokoll \lstinline[style=fileinline]|IPPROTO_DTLS_1_2| angegeben, sodass sich der Socket nach au{\ss}en wie ein normaler Datagramm-Socket benutzen l{\"a}sst (\lstinline[style=fileinline]|bind()|, \lstinline[style=fileinline]|recvfrom()|, \lstinline[style=fileinline]|sendto()|). Der TLS-Handshake erfolgt jedoch intern {\"u}ber mbedTLS. Bevor dieser Socket jedoch angelegt wird, wird ein Pre-Shared Key und eine dazugeh{\"o}rige PSK-Identit{\"a}t {\"u}ber die Funktion \lstinline[style=fileinline]|tls_credential_add()| im TLS-Credential-Subsystem registriert. Anschlie{\ss}end wird es dem Socket dann {\"u}ber \lstinline[style=fileinline]|setsockopt()| mit dem Optionsnamen \lstinline[style=fileinline]|TLS_SEC_TAG_LIST| zugewiesen. Dar{\"u}ber hinaus wird der Socket {\"u}ber die Option \lstinline[style=fileinline]|TLS_DTLS_ROLE| als Server konfiguriert. Der dabei benutzte Pre-Shared Key ist hierbei ein fest hinterlegter Demo-Schl{\"u}ssel im Quellcode, da es sich hierbei nur um Testzwecke handelt. In einer echten Anwendung w{\"u}rde ein solcher Schl{\"u}ssel bei einer Ger{\"a}teinbetriebnahme (Provisioning) individuell erzeugt werden.

\subsection{Firmware-Update {\"u}ber Sysbuild, MCUboot und MCUmgr}
Damit diese Anwendung, wie in Abschnitt 2.7.4 beschrieben, auch Firmware-Updates unterst{\"u}tzt, wird sie {\"u}ber das Sysbuild-System gebaut. Dieses erzeugt aus dem eigentlichen Anwendungscode und einem weiteren MCUboot-Bootloader-Image ein gemeinsames Firmware-Abbild. Das Anwendungs-Image selbst wird hierbei daf{\"u}r passend f{\"u}r die Flash-Partitionen gebaut und signiert, die von MCUboot verwaltet werden. Auf der Anwendungsseite wird zudem das MCUmgr-Management-Subsystem eingebunden, darunter die Bildverwaltungs- und die Betriebssystem-Kommandogruppe. Dar{\"u}ber kann ein neues Firmware-Abbild empfangen und anschlie{\ss}end ein Neustart ausgel{\"o}st werden. F{\"u}r dieses Management-Protokoll wird hierbei als {\"U}bertragungsweg UDP benutzt. Dadurch kann ein Firmware-Update {\"u}ber denselben Thread-Netzwerkpfad {\"u}bertragen werden, {\"u}ber den auch der CoAP-Sensordienst erreichbar ist.
